\section{从 Gaussian 到卡方与 Student \(t\)}
\label{sec:05-Probability/06-Common-Distributions/06}

你刚接手一批测量数据：\(2.1, 1.8, 2.3, 1.9, 2.0\)，五个数。算个均值 \(\bar x=2.02\)，很顺手。现在领导问你：真实均值可能在什么范围内？你翻出 Gaussian 置信区间的配方：\(\bar X\pm z_{\alpha/2}\,\sigma/\sqrt{n}\)。\(\bar X\) 有了，\(n=5\)，\(z_{0.025}=1.96\)——等等，\(\sigma\) 是多少？

你并不知道总体的标准差。你手上只有这五个数。唯一的办法是用样本标准差 \(s\) 去顶替 \(\sigma\)。\(s\) 算出来大约 \(0.21\)。于是你写下

\[
\frac{\bar X-\mu}{s/\sqrt{n}},
\]

然后下意识地把它当标准正态查表。这个下意识的动作需要停一下。

把 \(\sigma\) 换成 \(s\)，分母就不再是一个常数了——\(s\) 本身就是随机变量，它从同一批数据里估计出来，和分子 \(\bar X\) 共享同一组样本的随机波动。你手里的这个比值，分布不再是标准正态。要搞清楚它到底是什么分布，我们得先理解 \(s^2\) 从 Gaussian 数据里长出来的时候，服从什么规律。

\subsection{从标准正态的平方到卡方}

先看一个更简单的问题。假设你有一个服从 \(\mathcal N(0,1)\) 的随机数生成器，独立抽了 \(\nu\) 个值 \(Z_1,\ldots,Z_\nu\)。把每个数平方再求和：

\[
Q=\sum_{i=1}^{\nu}Z_i^2.
\]

这个 \(Q\) 能告诉我们什么？它的期望很直白：\(\E[Z_i^2]=\Var(Z_i)=1\)，所以 \(\E[Q]=\nu\)。但它的完整分布有一个名字：自由度 \(\nu\) 的卡方分布，记作

\[
Q\sim\chi^2_\nu.
\]

卡方是 Gamma 分布家族成员——取形状参数 \(\nu/2\)，速率参数 \(1/2\) 的那一条：

\[
\chi^2_\nu = \mathrm{Gamma}\Bigl(\frac\nu2,\;\text{rate }\frac12\Bigr).
\]

从这个血缘关系立刻读出 \(\Var(Q)=\nu\cdot(1/2)/(1/2)^2=2\nu\)。自由度越高，分布的中心越向右移，散布也越宽——但相对散布 \(\sqrt{\Var(Q)}/\E[Q]=\sqrt{2/\nu}\) 反而在收窄。

\subsection{自由度是什么意思}

``自由度 \(\nu\)'' 听起来抽象，但有一个几何直觉。\(Z_1,\ldots,Z_\nu\) 是 \(\nu\) 维空间中的一个随机点，各坐标独立标准正态。\(Q\) 是这个点到原点距离的平方。自由度的字面意思：有 \(\nu\) 个坐标可以独立变动，每个贡献一个 \(Z_i^2\)，合计 \(\nu\) 个独立方向。

如果这些坐标被某个线性约束绑住——比如要求 \(\sum_i Z_i=0\)——那么实际上只有 \(\nu-1\) 个方向是自由的。平方和的分布就从 \(\chi^2_\nu\) 降为 \(\chi^2_{\nu-1}\)。每一条独立的线性约束，吃掉一个自由度。

\subsection{样本方差的自由度：为什么是 \(n-1\)}

回到那五个数。设 \(X_1,\ldots,X_n\) 独立同分布于 \(\mathcal N(\mu,\sigma^2)\)。你算出

\[
\bar X=\frac1n\sum_i X_i,\qquad
S^2=\frac1{n-1}\sum_i (X_i-\bar X)^2.
\]

\(S^2\) 的分子里是 \(n\) 个残差 \(X_i-\bar X\) 的平方和。这些残差受一条约束：

\[
\sum_{i=1}^n (X_i-\bar X)=0.
\]

\(n\) 个残差，只有 \(n-1\) 个可以自由变化——给定任意 \(n-1\) 个，最后一个由和为零唯一确定。这就是 \(S^2\) 的自由度是 \(n-1\) 的几何来源。

严谨结论需要正交投影或 Cochran 定理，但核心事实可以直接陈述：对 Gaussian 样本，

\[
\frac{(n-1)S^2}{\sigma^2} \sim \chi^2_{n-1}.
\]

标准化后的样本方差，精确服从自由度为 \(n-1\) 的卡方分布。注意这个结论依赖 Gaussian 假设——是正态数据的特殊结构让残差平方和恰好落在卡方族里，不是随便什么分布都能享受的。

还有一个更惊人的事实：对 Gaussian 样本，\(\bar X\) 和 \(S^2\) 是独立的随机变量。

\[
\bar X \perp S^2.
\]

均值估计和方差估计互不干扰。这个性质在整个统计推断里价值连城——它意味着你可以分别讨论位置和尺度，两个估计之间没有信息纠缠。但它同样是 Gaussian 专属的奢侈品。

\subsection{Student \(t\)：把未知 \(\sigma\) 换掉的代价}

现在回到最初的问题。你要标准化 \(\bar X\)，但不知道 \(\sigma\)。很自然的操作是用 \(S\) 替代 \(\sigma\)：

\[
\frac{\bar X-\mu}{\sigma/\sqrt{n}}
= \frac{Z}{\text{常数}}
\sim \mathcal N(0,1),
\]

但

\[
\frac{\bar X-\mu}{S/\sqrt{n}}
= \frac{Z}{\sqrt{V/\nu}}
\]

不再是标准正态。分子是 \(Z\sim\mathcal N(0,1)\)，分母里 \(V\sim\chi^2_\nu\)，且二者独立（感谢 \(\bar X\perp S^2\)）。这个比率定义的分布叫自由度为 \(\nu\) 的 Student-\(t\) 分布：

\[
T = \frac{Z}{\sqrt{V/\nu}} \sim t_\nu.
\]

对 Gaussian 样本，

\[
T = \frac{\bar X-\mu}{S/\sqrt{n}} \sim t_{n-1}.
\]

现在可以回答那个关键问题：把 \(t\) 分布当标准正态用，会错在哪里？

\subsection{尾部：\(t\) 比 Gaussian 厚多少}

\(t\) 分布中心在零，对称，钟形——跟标准正态长得像。差别在尾部。\(S\) 在分母里随机波动：偶尔它偏小，低估了真实尺度 \(\sigma\)，于是比值 \(|\bar X-\mu|/(S/\sqrt{n})\) 就会被放大。这种放大效应让 \(t\) 分布的极端值比标准正态更常见。

拿具体数字来感受。自由度 \(\nu=5\) 的 \(t\) 分布，\(P(|T|>2)\) 大约 \(0.10\)；标准正态同等情形大约 \(0.05\)——相差一倍。如果你用 \(z_{0.025}=1.96\) 去构造 \(95\%\) 置信区间，实际覆盖率不足 \(90\%\)。操作上就是把本不该拒绝的假设错误地拒绝了，假阳性翻倍。自由度越小，这个差距越大。

再看另一头：\(\nu=30\) 时 \(P(|T|>2)\approx0.054\)，已经靠 Gaussian 很近了。自由度大到一定程度，\(S\) 对 \(\sigma\) 的估计足够稳定，分母几乎就是常数，\(t\) 和正态的差异在实用意义上消失（收敛理论见\hyperref[sec:05-Probability/08-Limit-Theorems/01]{``随机变量的多种收敛''}）。

\subsection{\(t\) 分布的矩也有存在条件}

一个分布能写出密度公式，不等同于它的矩都会老实存在。\(t_\nu\) 的尾部和 Cauchy 同源（\(\nu=1\) 就是 Cauchy），所以矩的存在性被自由度卡住：

\begin{itemize}
\tightlist
\item \(\nu\le 1\)：均值不存在——积分 \(\E[|T|]\) 发散；
\item \(1<\nu\le 2\)：均值存在（为零），但方差无穷；
\item \(\nu>2\)：方差才有限，\(\Var(T)=\dfrac{\nu}{\nu-2}\)。
\end{itemize}

小样本用 \(t\) 检验的时候，不仅要检查正态性，自由度过低时连均值的可估计性都有问题——只不过实践中 \(\nu\) 通常至少是 \(n-1\) 且 \(n\) 不会小到 2，这个坑在日常 \(t\) 检验里不常踩中，但它的存在本身就是一堂课：分布的存在和矩的存在是两笔独立的账。

\subsection{统计应用与模型条件}

卡方有三类经典应用：Gaussian 方差推断、拟合优度检验、二次型分布的推导。Student-\(t\) 的核心场景是未知方差下的均值区间估计和回归系数检验。当涉及多个方差比较时，两个独立卡方之比的分布导出 \(F\) 分布——它是 ANOVA 和模型选择检验的骨架。

所有这些精确结论都踩着同一个前提：独立同分布的 Gaussian 样本。大样本时，中心极限定理让部分方法对轻微偏离保持稳健——\(t\) 检验在对称分布、适度样本量下表现尚可。但大样本不是万能药：重尾会让样本方差本身拒绝稳定，强偏斜会让均值和方差耦合，依赖和异方差会破坏 \(\bar X\perp S^2\)。统计软件里一个 \(p=0.032\) 的 \(t\) 检验结果，背后始终悬着一串模型条件，点击``运行''之前值得逐条过一遍。

\subsection{分布族不是一份孤立的名单}

回看整个常见分布单元，把出生的血缘关系串起来，比单独背每个分布的密度有用得多：

\begin{itemize}
\tightlist
\item Bernoulli 的和产生 Binomial——每个试验扔一次硬币，累计正面数；
\item 大 \(n\)、小 \(p\) 的 Binomial 极限产生 Poisson——稀有事件的计数近似；
\item Poisson 过程的首次等待时间产生 Exponential，第 \(k\) 次等待产生 Gamma——时间和计数是一个硬币的两面；
\item 标准 Gaussian 平方和产生卡方——正态数据的尺度信息藏在平方里；
\item Gaussian 除以独立卡方尺度产生 Student \(t\)——用数据自己的波动去标准化自己；
\item 独立 Gamma 的归一化产生 Beta 或 Dirichlet——比例的分布。
\end{itemize}

这些联系让参数、矩和适用场景互为印证，而不是三组需要分别记忆的孤立知识点。真正掌握一个分布，是能从它的生成机制重新推出来——比如看见``独立 Gamma 归一化''就知道 Beta 的密度应该长什么样——而不是只认得名字和公式。

本篇的起点是五个数求均值的日常操作，终点是一个比表面看起来更深的结论：一旦你用同一批数据同时估计均值和方差，你就不再生活在 Gaussian 世界里。Student \(t\) 替你付了这笔用 \(S\) 替代 \(\sigma\) 的代价，但前提是数据确实来自正态分布。下一组篇目会追问：如果连这个前提也不成立呢？
